\section{Agent identity and delegated authority}
\label{sec:agent-identity}

Two different things are called an \emph{agent} in this paper, and the difference decides what each
one needs. The first is the human-facing assistant of \S19: it turns a sentence into an \appspec,
renders it for review, and hands the signing to the person's own device. That assistant holds no key
and needs no delegated authority, because the human's signature \emph{is} the authorisation; its open
problem is reviewability, not delegation. The second is autonomous software — a process that
provisions capacity, pays for it, renews it and retires it on its own schedule, with no person in the
loop at any step, on timescales where waiting for a human tap is itself the failure. Only the second
needs a delegation primitive. This section is about the second, exclusively, and the two must not be
allowed to blur: a property proved about a keyless assistant says nothing about a process that holds
a signing key, and the deployed system already contains one of each.

The section opens on a statement about what is deployed, and it is not a comfortable one.
\textbf{Authority in the deployed system is bounded by action and unbounded by value.} There exists,
today, a real scoped and revocable delegation of \emph{actions}; there exists no bound whatsoever on
\emph{value}; and there is already in production a bearer credential that confers unbounded signing
authority over a user's Flux identity. What follows is therefore a specification for a capability
that already ships without one, not a description of a mechanism the network has.

\subsection{What exists today}
\label{sec:agent-identity:today}

\Cref{tab:agent-today} is the baseline. Every row is \shipped and carries code provenance; the one
row that records an absence carries the search that established it.

\begin{table}[H]
\centering\footnotesize
\begin{tabular}{@{}p{0.235\textwidth} p{0.72\textwidth}@{}}
\toprule
\textbf{Capability} & \textbf{What is deployed} \\
\midrule
Human key custody, UTXO &
\shipped 2-of-2 ECDSA over secp256k1 in a sorted \texttt{OP\_CHECKMULTISIG} script. The wallet
builds and partially signs; the Key device signs the remaining input, finalises, and broadcasts
directly to a chain backend. The relay never sees a complete UTXO transaction.
\srccode{ssp-wallet/\allowbreak{}src/\allowbreak{}lib/\allowbreak{}constructTx.ts}{242--298},
\srccode{ssp-key/\allowbreak{}src/\allowbreak{}lib/\allowbreak{}constructTx.ts}{157--169} \\[2pt]

Human key custody, EVM &
\shipped Two-round MuSig-style Schnorr aggregation, verified on chain inside an ERC-4337 smart
account through the \texttt{ecrecover} trick, at an address made identical on all seven chains of the
deployment manifest by \texttt{CREATE2}.
\srccode{account-abstraction/\allowbreak{}contracts/\allowbreak{}schnorr/\allowbreak{}Schnorr.sol}{7--26},
\srccode{account-abstraction/\allowbreak{}contracts/\allowbreak{}MultiSigSmartAccount.sol}{145--159},
\srcdoc{account-abstraction/deployments.md:5--13} \\[2pt]

Human key custody, Solana &
\shipped 2-of-2 Ed25519 over a program-derived vault. Registration is permissionless — the vault
address is a PDA over $(\mathrm{sorted\ members}, \mathrm{threshold})$ — and the threshold is checked
only at execution.
\srccode{Solana-Multisig/\allowbreak{}programs/\allowbreak{}solana-multisig/\allowbreak{}src/\allowbreak{}lib.rs}{89--152,468--471},
\srccode{ssp-wallet/\allowbreak{}src/\allowbreak{}lib/\allowbreak{}wallet.ts}{489} \\[2pt]

Relay capability &
\shipped The relay holds no private key and cannot sign on any chain. It can read, store and
\emph{withhold} every payload it forwards. Several coordination endpoints and both socket
\texttt{join} paths still accept unauthenticated requests, marked in the shipped source as
\texttt{TODO:\allowbreak{} Make auth required after clients are updated}.
\srccode{ssp-relay/\allowbreak{}src/\allowbreak{}lib/\allowbreak{}socket.ts}{31--38},
\srccode{ssp-relay/\allowbreak{}src/\allowbreak{}routes.ts}{253--301} \\[2pt]

\textbf{Action} delegation &
\shipped The \FluxEdge API issues bearer tokens whose privileges are enumerated per key
(\texttt{start\_rental}, \texttt{stop\_rental}, \texttt{start\_deployment}, \texttt{stop\_deployment},
\texttt{get\_balance}) and enforced per route; \SI{2}{\per\second} per key and per IP; five
rate-limit violations inside \SI{5}{\minute} revoke the key automatically and notify its owner.
\srccode{pouwbackend/\allowbreak{}rest\_api/\allowbreak{}ratelimit.go}{41--45,203--222,375--390} \\[2pt]

\textbf{Value} delegation &
\emph{Absent.} No session key, spending cap, allowance, time-lock, guardian or social-recovery
construct exists in the wallet, key, relay, account-abstraction or Solana-multisig code.
\srccode{ssp-wallet, ssp-key, ssp-relay, account-abstraction, Solana-Multisig}{targeted search, no
match} \\[2pt]

Policy engine &
\shipped The enterprise M-of-N vault has priority-ordered rules with actions
\texttt{allow}\slash\texttt{block}\slash\texttt{require\_approval}\slash\texttt{timelock}, schedule windows,
rolling 24\,h/7\,d/30\,d transaction-count velocity limits, whitelists and named approval groups. It
decides whether a \emph{proposal} needs approval; it never signs, and every allowed proposal still
requires the vault's device signatures.
\srccode{ssp-relay-enterprise/\allowbreak{}src/\allowbreak{}services/\allowbreak{}policyEngineService.ts}{9--45},
\srccode{ssp-relay/\allowbreak{}src/\allowbreak{}routes.ts}{644--649} \\[2pt]

Gas sponsorship hook &
\shipped \texttt{IPaymaster} and \texttt{paymaster\allowbreak{}And\allowbreak{}Data} are fully wired through the account; no
contract in the repository implements a paymaster.
\srccode{account-abstraction/\allowbreak{}contracts/\allowbreak{}erc4337/\allowbreak{}interfaces/\allowbreak{}IPaymaster.sol}{1--49} \\[2pt]

De-facto delegation, in production &
\shipped A Firebase bearer token authorises remote signing at
\texttt{service.\allowbreak{}fluxcore.\allowbreak{}ai/\allowbreak{}api/\allowbreak{}signMessage}, which signs on the user's behalf. The shipped MCP tools
\texttt{deploy\_app}, \texttt{update\_app} and \texttt{renew\_app} use exactly this path; their own
description is ``Revalidate, Firebase-sign, register''.
\srccode{deploy-with-git-frontend/\allowbreak{}server/\allowbreak{}flux/\allowbreak{}fluxSession.js}{60--73},
\srccode{deploy-with-git-frontend/\allowbreak{}server/\allowbreak{}mcp/\allowbreak{}createServer.js}{174--266} \\
\bottomrule
\end{tabular}
\caption{Authority primitives in the deployed system. \shipped rows carry code provenance; the
absence row carries the search that established it.}
\label{tab:agent-today}
\end{table}

Three notes on that table, because each of them constrains what the rest of the section may claim.

First, the EVM contracts are reported audited by Halborn with the engagement finalised in February
2025, with findings in code since removed \srcdoc{account-abstraction/aa-schnorr-multisig-sdk/README.md:76--83};
that is documentation provenance and is not independently verified here. In-repo deployment
artifacts were observed for four of the manifest's networks, the fifth directory being a local
\texttt{hardhat} network
\srccode{account-abstraction/\allowbreak{}deployments/\allowbreak{}}{directory listing}, so ``identical on seven chains'' is a
manifest claim, not an artifact count.

Second, the \FluxEdge token model is a genuine bounded, revocable delegation, and the section must
not pretend otherwise. It has an enumerated privilege set, a per-key rate limit, and automatic
revocation on abuse with the owner notified. Two gaps keep it from being sufficient for an
autonomous agent: key issuance is not part of the documented REST surface, so a human must mint a key
through the web dashboard before an agent can use the API at all
\srcdoc{pouwbackend/rest\_api/v1/openapi.yaml:61}; and the privileges it scopes are actions on a
custodial credit balance, not signatures on a chain.

Third — and this is the row the whole section exists because of — the deployed delegation on the
\FluxCore path is unbounded. A holder of a valid Firebase bearer token can construct, sign and submit
a registration or update message end to end, with no client-side wallet interaction. There is
\textbf{no per-transaction cap, no rate limit, no cumulative budget, no scope over what may be
deployed, and no revocation short of invalidating the token}. The bound on a compromised token today
is whatever the underlying identity can authorise, which is everything that identity can authorise.
The local-signing paths (\SSP\ extension, ZelCore deep link) do not have this property, because the
key never leaves the user's device \srccode{deploy-with-git-frontend/\allowbreak{}src/\allowbreak{}services/\allowbreak{}deployService.js}{697--772};
the property is per-path, and \S19 states it per-path.

\subsection{The two authorities}
\label{sec:agent-identity:two}

The structural claim of this section is stated early because it organises everything after it. A
deployment involves \emph{two separable authorities}, and they are enforced in different places by
different parties, so they have different feasibility.

\begin{itemize}
  \item \textbf{Deployment authority (A)} --- what may run, at what resource vector, in how many
  instances, where, and under whose name. It is enforced by every \FluxOS node, at the moment the
  node validates the owner signature over an \appspec.
  \item \textbf{Spending authority (B)} --- how much \flux\ may move, to whom, and how fast. It is
  enforced by \Flux\ L1 consensus rules, which are those of a Zcash-derived UTXO chain.
\end{itemize}

\Cref{tab:two-authorities} sets them side by side. Conflating them is what makes agent delegation
look uniformly hard. Separated, (A) is tractable with the machinery that already runs and (B) is not
tractable at all on the \Flux\ L1 without either a new trusted party or new consensus state. The rest
of this section treats them separately and never averages them.

\begin{table}[H]
\centering\footnotesize
\begin{tabular}{@{}p{0.19\textwidth} p{0.36\textwidth} p{0.36\textwidth}@{}}
\toprule
 & \textbf{(A) Deployment authority} & \textbf{(B) Spending authority} \\
\midrule
Object bounded & the \appspec\ an agent may submit & the value an agent may move \\
Enforcement point & \FluxOS\ specification validation & \Flux\ L1 script validation \\
Enforcing party & every validating node, independently & consensus \\
Inputs available at enforcement & the message, the certificate, chain state at $\hgt$ & the spending
transaction and the outputs it consumes, only \\
Can enforce history-dependent limits? & yes --- a node sees the chain & no --- see
\cref{prop:budget-balance} \\
Feasibility class (\S\ref{sec:feasibility}) & Extension (row 12) & Research (row 13) \\
\bottomrule
\end{tabular}
\caption{The two authorities, their enforcement points, and why they have different feasibility.}
\label{tab:two-authorities}
\end{table}

\subsection{The delegation certificate}
\label{sec:agent-identity:cert}

Everything from here to \cref{sec:agent-identity:compare} is \planned. Nothing in it is deployed, no
part of it is scheduled to a height, and no claim below should be read in the present tense as a
statement about the network.

\begin{definition}[Delegation certificate]
\label{def:cert}
Let $\principal$ be the human principal, holding the 2-of-2 key set $(k_1,k_2)$ split across the
wallet and key devices of \cref{tab:agent-today}. Let $\agent$ be an autonomous agent holding a
single key pair $(sk_\agent, pk_\agent)$ generated on the agent's own host. A \emph{delegation
certificate} is the tuple
\[
  \mathsf{cert} \;=\; \bigl(pk_\agent,\; \scopeset,\; \budget,\; [t_0,t_1],\; \mathsf{rev},\; \nu\bigr),
  \qquad
  \sigma_{\mathsf{cert}} \;=\; \mathrm{Sign}_{k_1 \wedge k_2}\bigl(H(\mathsf{cert})\bigr),
\]
where
$\scopeset$ is the scope of permitted deployments (\cref{def:scope});
$\budget = (b_{\mathrm{tx}}, b_{\mathrm{rate}}, b_{\mathrm{total}}) \in
(\mathbb{R}_{\ge 0}\flux) \times (\mathbb{R}_{\ge 0}\flux\,\text{per block}) \times
(\mathbb{R}_{\ge 0}\flux)$ is the intended value bound, being respectively a per-transaction cap, a
rate limit and a cumulative budget;
$[t_0,t_1] \subset \mathbb{N}$ is a validity window in block heights with $t_0 < t_1$;
$\mathsf{rev}$ is a revocation anchor, an outpoint on the \Flux\ L1 controlled by $\principal$ and
unspent at $t_0$ (\cref{def:revoked});
and $\nu \in \mathbb{N}$ is a monotone certificate counter, so that re-issue to the same
$pk_\agent$ is unambiguous and an older certificate can never be presented as a newer one.
$H$ is the canonicalising hash of the \appspec\ fingerprint recipe (\S10) and
$\mathrm{Sign}_{k_1 \wedge k_2}$ denotes the full 2-of-2 ceremony, not a single-device signature.
\end{definition}

\begin{assumption}[Model, from \S2]
\label{ass:agent}
The specification of this section is stated against the following, each of which is fixed in the
system model of \S2 and not re-argued here.
\begin{enumerate}
  \item[(A1)] \emph{Principal custody.} Producing $\sigma_{\mathsf{cert}}$ requires both $k_1$ and
  $k_2$. Compromise of one device does not yield a certificate.
  \item[(A2)] \emph{Node validation.} Every \FluxOS\ node independently validates owner signatures
  over \appspec s and has read access to \fluxd\ chain state at the height it validates. This is the
  deployed behaviour described in \S9 and \S10.
  \item[(A3)] \emph{Chain state.} The \Flux\ L1 is a Zcash-derived UTXO chain with no covenant
  opcodes, no transaction introspection, and no account abstraction. Reorganisations are bounded by
  \texttt{MAX\_REORG\_LENGTH} $=40$ blocks
  \srccode{fluxd/\allowbreak{}src/\allowbreak{}main.cpp}{8983--8988}, with the historical override to $5{,}000$ over heights
  $[2{,}020{,}000,\,2{,}025{,}000]$ noted in \S5.
  \item[(A4)] \emph{No escrow.} $sk_\agent$ is generated on the agent's host and is never
  transmitted. Any variant in which a service can sign for $\agent$ is a custody design and is
  labelled one.
  \item[(A5)] \emph{Adversary.} The adversary $\adv$ of \S2 may compromise the agent host, may
  operate or coerce the relay, and is computationally bounded so that signature forgery is
  infeasible. $\adv$ does not hold $k_1$ and $k_2$ simultaneously.
\end{enumerate}
\end{assumption}

\subsection{Deployment authority}
\label{sec:agent-identity:deploy}

\planned. Deployment authority is the half of the problem that composes with machinery the network
already runs.

\begin{definition}[Scope and containment]
\label{def:scope}
The scope is the tuple
\[
  \scopeset \;=\; \bigl(\mathcal{N},\; \resvec^{\max},\; \ninst^{\max},\; \mathcal{I},\;
  \dur^{\max},\; \mathcal{G}\bigr),
\]
where $\mathcal{N}$ is a finite set of permitted application names;
$\resvec^{\max} \in \mathbb{R}^3_{>0}$ is a componentwise ceiling on the resource vector
$\resvec = (\text{cpu}, \text{memory}, \text{rootFs})$;
$\ninst^{\max} \in \mathbb{N}$ bounds the instance count $\ninst$;
$\mathcal{I}$ is a finite set of permitted container images, each identified by content digest and
not by tag;
$\dur^{\max} \in \mathbb{N}$ bounds the requested duration $\dur$ in blocks;
and $\mathcal{G}$ is a set of permitted geolocation constraints in \FluxOS's \texttt{acXX}\slash\texttt{a!cXX}
form.
A specification $s$ is \emph{contained} in $\scopeset$, written $s \in \scopeset$, iff its name lies
in $\mathcal{N}$, its resource vector satisfies $\resvec(s) \le \resvec^{\max}$ componentwise, its
instance count satisfies $\ninst(s) \le \ninst^{\max}$, every image digest named by $s$ lies in
$\mathcal{I}$, $\dur(s) \le \dur^{\max}$, and the geolocation constraints of $s$ lie in $\mathcal{G}$.
Containment is a predicate on $s$ and $\scopeset$ alone: it reads nothing else.
\end{definition}

\begin{definition}[Acceptance]
\label{def:accept}
A node evaluating an agent-submitted specification $s$, carrying an agent signature $\sigma_s$, a
certificate $\mathsf{cert}$ and its principal signature $\sigma_{\mathsf{cert}}$, at height $\hgt$,
computes
\begin{multline*}
  \mathrm{Accept}(s,\hgt) \;\equiv\;
  \mathrm{Verify}_{pk_\agent}\bigl(H(s), \sigma_s\bigr) \;\wedge\;
  \mathrm{Verify}_{\mathsf{owner}}\bigl(H(\mathsf{cert}), \sigma_{\mathsf{cert}}\bigr) \\
  {}\wedge\; t_0 \le \hgt \le t_1 \;\wedge\;
  s \in \scopeset \;\wedge\;
  \neg\,\mathrm{Revoked}(\mathsf{cert},\hgt),
\end{multline*}
where $\mathsf{owner}$ is the owner identity already carried by the specification and
$\mathrm{Revoked}$ is \cref{def:revoked}.
\end{definition}

\begin{algorithm}[H]
\DontPrintSemicolon
\KwIn{message $m = (s, \sigma_s, \mathsf{cert}, \sigma_{\mathsf{cert}})$; local height $\hgt$;
local UTXO set $U(\hgt)$; owner public key set for $\mathsf{owner}$}
\KwOut{\textsc{accept} or \textsc{reject}}
\tcp{Invariant: the result is a deterministic function of $(m, \hgt, U(\hgt))$ alone.
No network fetch, no registry, no service.}
\lIf{$\neg\,\mathrm{Verify}_{\mathsf{owner}}(H(\mathsf{cert}), \sigma_{\mathsf{cert}})$}{\Return
\textsc{reject}}
\lIf{$\hgt < t_0 \vee \hgt > t_1$}{\Return \textsc{reject}}
\tcp{revoked, or anchor never existed --- both fail closed}
\lIf{$\mathsf{rev} \notin U(\hgt)$}{\Return \textsc{reject}}
\lIf{$\neg\,\mathrm{Verify}_{pk_\agent}(H(s), \sigma_s)$}{\Return \textsc{reject}}
\lIf{$\mathsf{owner}(s) \neq \mathsf{owner}(\mathsf{cert})$}{\Return \textsc{reject}}
\ForEach{clause $c$ of \cref{def:scope}}{
  \lIf{$s$ violates $c$}{\Return \textsc{reject}}
}
\Return \textsc{accept}\;
\caption{\textsc{AcceptDelegatedSpec}. \textbf{Failure mode:} every branch fails closed. A node that
cannot evaluate $\mathsf{rev} \in U(\hgt)$ --- for instance one whose local state does not reach the
anchor's height --- rejects rather than accepts, so a partially-synced node under-accepts and never
over-accepts. \planned}
\label{alg:accept}
\end{algorithm}

\begin{property}[Local enforceability]
\label{prop:local}
$\mathrm{Accept}(s,\hgt)$ is computable by any validating node from the submitted message, the
node's own chain state at $\hgt$, and nothing else. It requires no consensus change and introduces
no trusted party beyond those \S2 already names.
\end{property}

\begin{proof}
Take the conjuncts of \cref{def:accept} in turn. The two $\mathrm{Verify}$ terms are signature checks
over data carried in the message, against $pk_\agent$ (carried in $\mathsf{cert}$, which is itself
signed) and against the owner key set the node already resolves in order to validate any \appspec\
(A2). The window test compares $\hgt$, which the node holds by definition, against two integers in
the message. Containment $s \in \scopeset$ is by \cref{def:scope} a predicate on $s$ and $\scopeset$
alone. $\mathrm{Revoked}$ is by \cref{lem:utxo-rev} a membership test in the node's own UTXO set at
$\hgt$. Every input is therefore either in the message or in chain state the node already maintains,
which establishes the first claim. For the second: the check is performed at the same point, by the
same code path, and on the same inputs as the owner-signature validation \FluxOS\ already performs
(A2), so it changes no consensus rule and adds no party --- the only new authority is $\principal$'s
own key set, which already authorises the specification today. \qedhere
\end{proof}

Deployment authority would fit \specv{9} cleanly: the certificate is one additional specification
field, validated where it is enforced, in a version that is already adding fields (\S10). That is the
recommended placement, and it is why \S\ref{sec:feasibility} classes this as Extension (row 12) rather than Engineering.

\begin{remark}[Digest pinning is necessary, and against a generic runner it is not sufficient]
\label{rem:digest}
$\mathcal{I}$ is defined over digests rather than tags for a concrete reason. The deployed tooling
emits \texttt{runonflux/\allowbreak{}orbit:latest}, an unpinned mutable tag, into every generated specification
\srccode{deploy-with-git-frontend/\allowbreak{}src/\allowbreak{}services/\allowbreak{}deployService.js}{462,453--490}. A scope that permits
a mutable tag permits whatever that tag resolves to at node pull time, which is to say arbitrary
code, and the delegation is then vacuous: the acceptance predicate would be satisfied by a
specification whose behaviour the principal never bounded.

Pinning the digest closes that hole and opens a second one, which the specification must name rather
than hide. The pinned image in question is a \emph{generic runner}: the Orbit container clones
whatever repository \texttt{GIT\_REPO\_URL} names and builds and runs it
\srccode{orbit/\allowbreak{}entrypoint.sh}{1--9,29}. Against such an image, fixing the digest fixes the
\emph{runtime} and not the \emph{workload}, so an agent whose scope permits one pinned generic-runner
digest can still cause arbitrary code to execute by varying an environment parameter. Bounding that
requires $\scopeset$ to constrain the environment-parameter set as well as the image set, and the
scope tuple of \cref{def:scope} does not. The requirement is stated; the construction is not.
\proposednote{The scope tuple gains a seventh component: an allowed-key set for the
environment-parameter map, with values unconstrained. Rationale: the keys are what a node can check
cheaply and deterministically at deployment, and they are what carries the risk --- an agent that can
set an arbitrary environment key can often reach configuration the owner never intended. Constraining
values as well would require the verifier to understand each application's semantics, which no node
can do. \srcinferred{} --- proposed by this paper on the same basis as \Cref{rem:cert-location}.}
\end{remark}

\subsection{Revocation}
\label{sec:agent-identity:revoke}

\begin{definition}[Revocation anchor]
\label{def:revoked}
$\mathsf{rev}$ is an outpoint on the \Flux\ L1 spendable by $\principal$ and required to be unspent
at $t_0$. The certificate is revoked at the first height at which that outpoint is spent:
\[
  \mathrm{Revoked}(\mathsf{cert},\hgt) \;\equiv\;
  \exists\, \hgt' \le \hgt \;:\; \mathsf{rev} \in \mathrm{Spent}(\hgt').
\]
\end{definition}

\begin{lemma}[UTXO-set formulation]
\label{lem:utxo-rev}
Let $U(\hgt)$ be the unspent-output set at height $\hgt$. Given the issuance precondition that
$\mathsf{rev} \in U(t_0)$, and given that a spent outpoint cannot be recreated,
\[
  \mathrm{Revoked}(\mathsf{cert},\hgt) \;\equiv\; \mathsf{rev} \notin U(\hgt)
  \qquad \text{for all } \hgt \ge t_0 .
\]
\end{lemma}

\begin{proof}
If $\mathsf{rev}$ is spent at some $\hgt' \le \hgt$ then it leaves $U$ at $\hgt'$ and, by
non-recreatability, does not re-enter, so $\mathsf{rev} \notin U(\hgt)$. Conversely, if
$\mathsf{rev} \notin U(\hgt)$ then, since $\mathsf{rev} \in U(t_0)$ and $t_0 \le \hgt$, the only way
to leave $U$ between $t_0$ and $\hgt$ is to be consumed by a transaction, i.e. spent at some
$\hgt' \le \hgt$. \qedhere
\end{proof}

The lemma matters for implementation rather than for theory: it turns revocation into a UTXO-set
membership test, which is exactly the state a pruned node retains (\S8), so revocation checking does
not force any node to keep history it would otherwise discard. A node that finds
$\mathsf{rev} \notin U(\hgt)$ without being able to distinguish ``spent'' from ``never existed''
rejects, which is the fail-closed direction.
\begin{remark}[Outpoint non-recreatability --- verified]
\label{rem:outpoint-verified}
\Cref{lem:utxo-rev} depends on a spent collateral outpoint not being recreatable, which was
previously assumed from inherited Bitcoin behaviour rather than checked. \fluxd{} enforces BIP-30
unconditionally in \texttt{ConnectBlock}: a block containing a transaction whose txid already exists
with unspent outputs is rejected as \texttt{bad-txns-BIP30}, with no height gate and no activation
flag \shipped \srccode{fluxd/src/main.cpp}{3868--3875}.

One nuance is worth stating precisely rather than glossing, because the rule permits overwrite when
the prior coins are \emph{completely} spent. For a non-coinbase outpoint, recreating the same txid
requires replaying a transaction with identical inputs, and those
inputs are by construction already spent, so the recreation cannot be built. The lemma therefore
holds for non-coinbase anchors on the deployed rules.
\end{remark}

\begin{property}[Monotone revocation, and its latency]
\label{prop:monotone}
Once $\mathrm{Revoked}(\mathsf{cert},\hgt)$ holds it holds at every $\hgt' > \hgt$, except under a
reorganisation deeper than $\hgt' - \hgt$. Revocation takes effect after one confirmation of the
spending transaction plus specification-propagation time.
\end{property}

\begin{proof}
Monotonicity is immediate from \cref{def:revoked}: the existential quantifier ranges over a set that
only grows with $\hgt$. The exception is equally immediate: if a reorganisation removes the block at
$\hgt'$ containing the spend, the existential no longer witnesses, and the certificate is live again
on the new chain. By (A3) this is bounded by \texttt{MAX\_REORG\_LENGTH} $=40$ blocks in normal
operation. Latency: the predicate becomes true for a node at the height at which that node accepts
the block containing the spend, so one confirmation, plus the time for a specification submitted
before that height to stop propagating --- during which a node that has not yet seen the spending
block will still accept. \qedhere
\end{proof}

At the tooling's figure of $2{,}880$ blocks per day
\srccode{deploy-with-git-frontend/\allowbreak{}src/\allowbreak{}services/\allowbreak{}deployService.js}{105}, one confirmation is
approximately \SI{30}{\second} \srcinferred. The honest statement of the guarantee is therefore
\emph{revocation is effective within one confirmation, and is reversible for up to $40$ blocks
thereafter}, and a principal revoking in response to a live compromise should assume the agent can
act for at least that long.

Two weaker designs were considered and rejected. \emph{Expiry alone} --- relying on $t_1$ and issuing
short certificates --- cannot stop an in-progress compromise: between detection and $t_1$ the agent
retains full scope, and shortening $t_1$ to make that window small multiplies re-issue traffic
against a 2-of-2 ceremony that requires two human devices. \emph{A revocation list served over
HTTPS} is operationally trivial and reintroduces a Flux-operated dependency in the validation path of
every node, which is precisely the class of dependency \S12 already inventories too many of; it would
also make revocation checking fail-open on service outage unless every node hard-fails on an
unreachable list, which is worse.

\subsection{Spending authority, and the honest result}
\label{sec:agent-identity:spend}

\begin{proposition}[Budget reduces to balance]
\label{prop:budget-balance}
On the \Flux\ L1 as deployed, delegating spending authority to a key not held by $\principal$ bounds
$\agent$'s spending by exactly the balance of the outputs that key can spend, and by nothing else.
Neither $b_{\mathrm{rate}}$ nor $b_{\mathrm{total}}$ of \cref{def:cert} is enforceable by consensus.
\end{proposition}

\begin{proof}[Proof sketch]
Enforcing $b_{\mathrm{rate}}$ or $b_{\mathrm{total}}$ requires a validation predicate that is a
function of the agent's \emph{transaction history} at spend time: whether the agent has already spent
$x$ in the preceding $n$ blocks, or $y$ in total since $t_0$. Script validation on a Bitcoin-derived
chain is a pure function of the spending transaction and the outputs it consumes; a script sees its
own input, its own transaction, and nothing about which other outputs the same key has spent. By
(A3), \Flux\ inherits this via Zcash with no covenant opcode, no introspection opcode and no account
abstraction, so no script can express the predicate, and no consensus rule in \fluxd\ supplies it
outside script. What remains enforceable is the trivial bound: a key can spend at most the outputs it
can spend. Hence the only real bound is the funded balance, which is $b_{\mathrm{tx}}$ collapsed into
$b_{\mathrm{total}}$ and both collapsed into the balance. \qedhere
\end{proof}

\Cref{prop:budget-balance} is the reason \S\ref{sec:feasibility} classes bounded agent spending as Research (row 13) and
not as an increment. The deployable construction is correspondingly modest, and is presented as such.

\begin{construction}[Funded allowance]
\label{con:allowance}
\planned. $\principal$ funds an address controlled solely by $sk_\agent$ with $b_{\mathrm{total}}$,
and tops it up on a schedule of the principal's choosing. The agent's maximum spending loss is that
balance. Revocation of spending authority is a sweep: $\principal$ spends the funded outputs back to
an address of their own. The sweep races any in-flight agent transaction, and the race is bounded by
one confirmation --- the same bound as \cref{prop:monotone}, and for the same reason.
\end{construction}

Three richer constructions exist. They are given as the design space, with costs, not as a
recommendation.

\begin{enumerate}
  \item \textbf{2-of-3 with a policy co-signer.} $\agent$ holds one key, $\principal$ a second, and a
  policy service the third; the service co-signs only transactions inside $\budget$. This is
  enforceable with what runs today, and the enterprise policy engine of \cref{tab:agent-today} is the
  obvious substrate --- it already evaluates velocity limits and schedule windows. \emph{Cost: a new
  trusted party.} The co-signer cannot steal, since it is one key of three, but it can refuse; and for
  an autonomous agent, refusal is a liveness failure that arrives precisely when the agent most needs
  to act. A design whose safety property is ``a third party can stop you'' has, for an unattended
  process, converted a theft risk into an availability risk rather than removing a risk.
  \item \textbf{Pre-signed rate-limited chain.} $\principal$ pre-signs a sequence of payments with
  increasing \texttt{nLockTime}, handing $\agent$ a spend schedule it cannot accelerate. No new trust
  is introduced and nothing new is required of consensus. \emph{Cost:} the amounts and the schedule
  are fixed at issue, which does not fit demand-driven provisioning --- an agent that needs
  \SI{40}{\percent} more capacity at hour three cannot buy it --- and revocation requires
  $\principal$ to broadcast a conflicting spend, which races the agent in the mempool.
  \item \textbf{Consensus-level delegation registry.} A new transaction type binding $pk_\agent$ to a
  budget held as consensus state, decremented as the agent spends. This is the clean answer: it makes
  $\budget$ real, it needs no third party, and it makes $b_{\mathrm{rate}}$ and $b_{\mathrm{total}}$
  the enforceable objects \cref{prop:budget-balance} says they cannot otherwise be. \emph{Cost:} a
  consensus change and new consensus state, on a chain that has just completed one major transition
  --- not cleanly (\S5) --- and has three further changes queued (\S8).
\end{enumerate}

The paper's position is: \cref{con:allowance} is the deployable baseline; (c) is the principled
endpoint; the choice between them is not made here.
\proposednote{\Cref{con:allowance} only --- the funded-allowance construction, with no co-signer,
pre-signed chain or consensus registry. This follows directly from \Cref{prop:budget-balance}: on this
chain the only enforceable bound is the balance of the key the agent holds, so the three richer
endpoints buy expressiveness the consensus layer cannot honour. A co-signer additionally reintroduces
a trusted third party, which is what the agent-native claim exists to remove
(\Cref{rem:agent-must-hold-keys}).}

\textbf{The asymmetry with the EVM path should be stated precisely rather than smoothed over.} On the
parallel-chain accounts, (B) is not hard: the ERC-4337 account already in production is a program
with persistent state, and it admits a session-key validator module and the sponsorship hook that
\texttt{IPaymaster} already declares. The account as deployed does not use it --- \texttt{\_validate\allowbreak{}Signature}
returns valid-with-no-time-bound or failure, never packing \texttt{validUntil}\slash\texttt{validAfter} into
\texttt{validationData} even though the base class supports it
\srccode{account-abstraction/\allowbreak{}contracts/\allowbreak{}MultiSigSmartAccount.sol}{145--159} --- but the substrate is
there. The \Flux\ L1's is not, and no amount of design closes that gap.

\subsection{Invariants}
\label{sec:agent-identity:invariants}

The specification is required to maintain the following. \planned throughout.

\begin{enumerate}
  \item \textbf{No escrow.} $sk_\agent$ is generated on the agent's host and never transmitted
  (A4). Any variant in which a service holds or can reconstruct $sk_\agent$ is a custody design and
  must be labelled one --- which is exactly what the \FluxCore\ signing path is, and is why
  \cref{tab:agent-today} names it.
  \item \textbf{Certificate authenticity.} Every accepted agent action traces to exactly one
  certificate signed by the full 2-of-2 principal key set. Single-device issuance is forbidden:
  a delegation issued more weakly than the custody it delegates from would make the delegation, not
  the custody, the attack surface.
  \item \textbf{Monotone revocation.} \Cref{prop:monotone}, with the reorg caveat stated and not
  hidden.
  \item \textbf{Scope containment.} $s \in \scopeset$ is evaluated by every validating node from the
  specification and the certificate alone, with no external lookup (\cref{prop:local}).
  \item \textbf{Loss bound.} \Cref{prop:loss} below.
\end{enumerate}

\begin{proposition}[Loss bound under full agent compromise]
\label{prop:loss}
Let $\adv$ fully compromise $\agent$ at height $\hgt_c \in [t_0,t_1]$, and let $\hgt_r \ge \hgt_c$ be
the height at which $\principal$'s revocation confirms. Under invariants 1--4 and
\cref{con:allowance}, the loss is bounded by
\[
  \mathrm{Loss}(\adv) \;\le\;
  \mathrm{bal}\bigl(sk_\agent, \hgt_c\bigr) \;+\;
  \price^{\max}\!\bigl(\scopeset,\, \min(t_1,\hgt_r) - \hgt_c\bigr),
\]
where $\mathrm{bal}(sk_\agent,\hgt_c)$ is the balance of outputs $sk_\agent$ can spend at
$\hgt_c$ and $\price^{\max}(\scopeset,\Delta)$ is the maximum cost of deployments admissible under
$\scopeset$ over $\Delta$ blocks. In words: \emph{the maximum loss from full compromise is the
balance the agent's key controls, plus the cost of any in-scope deployment for the remaining validity
window.}
\end{proposition}

\begin{proof}
By (A4) and invariant 1, $\adv$ obtains $sk_\agent$ and nothing else; by (A5) $\adv$ cannot produce
$\sigma_{\mathsf{cert}}$ for any certificate $\principal$ did not sign, and cannot forge
$\sigma_s$ for a different key. So $\adv$'s capabilities are exactly $\agent$'s. Spending: by
\cref{prop:budget-balance} and \cref{con:allowance}, $\adv$ can move at most the outputs
$sk_\agent$ controls, i.e. $\mathrm{bal}(sk_\agent,\hgt_c)$, plus any top-up $\principal$ makes after
$\hgt_c$ --- which the statement excludes by evaluating the balance at $\hgt_c$ and which a principal
responding to a detected compromise will not make. Deployment: by \cref{def:accept} every accepted
specification satisfies $s \in \scopeset$ and $t_0 \le \hgt \le t_1$ and
$\neg\mathrm{Revoked}$, so the admissible set is bounded by $\scopeset$ and the window closes at
$\min(t_1,\hgt_r)$ --- at $t_1$ by expiry, at $\hgt_r$ by \cref{prop:monotone}. Summing the two
disjoint capabilities gives the bound. \qedhere
\end{proof}

Two qualifications travel with \cref{prop:loss} and must not be dropped when it is quoted. First, the
deployment term is only as tight as $\scopeset$, and by \cref{rem:digest} a scope containing a
generic-runner image is not tight at all. Second, $\hgt_r$ is not under the principal's full control:
it depends on detection, on a 2-of-2 ceremony requiring two devices, and on one confirmation.

\subsection{Attack analysis}
\label{sec:agent-identity:attacks}

\paragraph{Compromised agent.} Covered by \cref{prop:loss}. The worst case is that $\adv$ spends the
funded balance and deploys anything inside $\scopeset$ until revocation confirms, each such deployment
then running to its own paid-through height (\Cref{rem:revoke-vs-lifecycle}). Digest-pinned
$\mathcal{I}$ is what stops ``anything inside $\scopeset$'' from meaning arbitrary code, subject to
\cref{rem:digest}. Revocation latency is bounded and stateable, which is the property that makes the
exposure quantifiable rather than open-ended.

\paragraph{Compromised assistant holding no key.} An assistant of the \S19 kind can propose a bad
specification and can do nothing else: it cannot deploy and it cannot spend, because the acceptance
predicate requires a signature it cannot produce. \textbf{This property holds only on the
local-signing path.} On the \FluxCore/Firebase path the assistant's bearer token \emph{is} signing
authority, and the property fails completely --- there is no key the assistant lacks. Any statement
of the form ``the assistant configures but never signs'' is a statement about one of two deployed
paths, and \S19 states it per path. This section's contribution to that discussion is narrow and
concrete: the delegated path is not a defect so much as an undocumented instance of exactly the
delegation specified here, running without any of the bounds specified here.

\paragraph{Malicious principal.} $\principal$ issues a certificate, lets $\agent$ act, then claims no
authorisation was given. The certificate is signed under (A1) and anchored on chain through
$\mathsf{rev}$, whose creation and spend are both public and heighted, so the claim is
non-repudiable up to compromise of both principal devices. No additional mechanism is required, and
none is specified.

\paragraph{Griefing the agent.} $\principal$ revokes mid-operation and strands a running deployment.
This is deliberate and is not treated as an attack: the principal must always be able to stop the
agent, and any design that lets an agent resist its principal is worse than the problem it solves.
Nothing happens to the \emph{workload} at that moment: it continues to its paid-through height under
the state machine of \S16, which is not restated here (\Cref{rem:revoke-vs-lifecycle}).
\begin{remark}[Reconciling revocation with \S16's lifecycle: they are orthogonal, and must stay so]
\label{rem:revoke-vs-lifecycle}
Revocation and the payment lifecycle both bear on a deployment's lifetime, so the case they must agree on is a
\textbf{revoked-but-paid} application: the owner has revoked the certificate under which the agent
deployed it, but the application's paid-through height has not been reached. Neither section
previously stated the transition, and the reconciliation is this: \textbf{the two state machines are
orthogonal; revocation gates acceptance, and \S\ref{sec:paying-for-capacity} alone governs a deployment once accepted.}

\S\ref{sec:paying-for-capacity}'s states $\{\textsf{funded},\textsf{grace},\textsf{suspended},\textsf{evicted}\}$
answer \emph{has this been paid for}. Revocation answers \emph{was this authorised}, and it is evaluated
once, by \cref{alg:accept} on a submitted message; nothing in this section re-evaluates it for a running
application. A specification is accepted only while both hold, so a revoked-but-paid application
\textbf{continues to its paid-through height $e_\app$ like any other paid application, and what
revocation stops is every further registration, update or renewal under that certificate}. It does not
pass through \textsf{grace} or \textsf{suspended} either: those two states exist to absorb transient
funding failures --- a bridge delay, a wallet outage --- and to give a paying tenant time to recover,
and nothing about the deployment's funding has changed. The cost of this narrowing is stated rather
than hidden: an owner who revokes a compromised agent's authority does not stop the deployments that
agent already had accepted; they run out their paid terms, and their cost is the $\price^{\max}$ term
of \cref{prop:loss}. Stopping them sooner would require every node to re-check
$\mathsf{rev} \in U(\hgt)$ for running applications at each height, a mechanism this paper does not
specify.

Two consequences follow, and both are stated rather than left implicit.
\begin{enumerate}
  \item \textbf{Nothing is forfeited on revocation.} Unlike moderation eviction
  (\Cref{rem:o16-burn}), revocation is the owner's own act against their own agent, not against the
  deployment: the paid term is consumed as for any other paid application, and there is no
  counterparty and no forfeiture event.
  \item \textbf{Revocation cannot be used to escape a moderation ban.} Because revocation touches
  acceptance and nothing else, an owner revoking their own certificate after a ban has been carried
  changes nothing: the application was already stopped by the ban, and revocation is not a path back
  to any state.
\end{enumerate}
\end{remark}

\paragraph{The relay as a withholding adversary.} The relay cannot sign (\cref{tab:agent-today}), so
it cannot forge a certificate. It can withhold, and for an autonomous agent withholding is a liveness
attack against certificate \emph{issuance}. The mitigation is structural rather than added:
certificates are issued once and used many times, and by \cref{prop:local} a node validating an
agent's specification consults no relay, so a withholding relay delays a new delegation and does not
interrupt a running one. Separately, and more urgently: the unauthenticated coordination endpoints
flagged in the shipped source \srccode{ssp-relay/\allowbreak{}src/\allowbreak{}routes.ts}{253--301} let a party who learns a
\texttt{wkIdentity} string join that identity's room and post payloads for it. Delegation should not
be built on top of those endpoints until they are closed, and this section treats closing them as a
precondition rather than as a parallel improvement.

\paragraph{Sybil agents.} Irrelevant by construction. Authority flows from a certificate, not from an
identity, so $\adv$ generating $n$ agent key pairs holds $n$ keys and zero authority; each would
require its own signature under (A1). This is the one adversary in \S2's catalogue that this
mechanism defeats for free, and the reason is worth naming: nothing in the design is per-identity
rate-limited or per-identity funded, so there is nothing for extra identities to multiply.

\subsection{Open parameters}
\label{sec:agent-identity:open}

\begin{remark}[Certificate encoding and location --- resolved by splitting the two authorities]
\label{rem:cert-location}
The three candidate locations --- a \specv{9} specification field, a separate broadcast message, or an
on-chain commitment with off-chain retrieval --- were carried as an undecided choice. They are
resolvable once the certificate's two halves are separated, because the halves have different
enforceability and therefore different homes.

\textbf{Deployment authority belongs in a \specv{9} specification field.} It is genuinely
enforceable: a node validating a deployment already holds the specification, so it can check the
scope tuple against what is being deployed at exactly the point where the constraint bites, with no
new message type, no second propagation path, and no retrieval availability problem. The cost is
real and bounded --- it binds the mechanism to the \specv{9} enforcement height (\S\ref{sec:appspec-v9})
--- and is worth paying, because the alternatives buy schedule independence with a permanent
consistency obligation between two propagation paths.

\textbf{Spending authority has no home, because it has no enforcement.}
\Cref{prop:budget-balance} shows that neither $b_{\mathrm{rate}}$ nor $b_{\mathrm{total}}$ is
enforceable by consensus on a Bitcoin-derived chain without covenants or introspection; the only
bound that survives is the balance of the key the agent can spend. Encoding an unenforceable bound
on-chain would state a guarantee the chain does not provide, which is worse than stating none. The
spending half of the certificate is therefore advisory: meaningful to an operator or a wallet that
chooses to honour it, and invisible to consensus.

This split is also what the team's own settled position implies. The protocol offers no automatic
renewal and no standing spend authority (\Cref{rem:no-protocol-renewal}); convenience of that kind
is a service-layer product. A certificate scheme that tried to carry enforceable spending limits
would be reintroducing at the consensus layer precisely what that decision removed.
\end{remark}

\begin{remark}[What this costs the agent-native claim, stated plainly]
\label{rem:agent-must-hold-keys}
Two settled decisions --- no protocol auto-renewal, and spending bounds that collapse to balance ---
have a consequence for the thesis that should be stated rather than left for a reader to derive. An
autonomous tenant that provisions, pays and renews without a human \textbf{must hold a signing key
and use it}, and the only way to bound its exposure is to fund that key with no more than the
operator is willing to lose. There is no protocol object that says ``this agent may spend at most
$x$ per day''; \Cref{prop:budget-balance} shows there cannot be one on this chain without a
consensus change of the class \S\ref{sec:feasibility} calls Research.

The alternative --- delegating to a custodial operator that renews on the tenant's behalf --- works
today and is what the team expects most users to choose \srcintent{design intake},
but it reintroduces exactly the intermediary the agent-native claim is about removing. The honest
statement of the target architecture is therefore narrower than ``agents transact without
intermediaries'': agents transact without intermediaries \emph{if they hold keys}, and the risk of
holding them is bounded only by funding discipline. The paper makes that claim and not the broader
one.
\end{remark}

\proposednote{Whether the scope $\scopeset$ also constrains where value may be sent: it does not.
\cref{def:scope} bounds what may be deployed and says nothing about
where value may go, and that asymmetry is kept deliberately. A destination constraint would be
unenforceable for the same reason a spending cap is (\Cref{prop:budget-balance}): script validation
cannot observe where a future transaction sends value. Writing an unenforceable constraint into the
scope tuple would state a guarantee the chain does not provide, which is worse than stating none.}

\proposednote{$t_1 - t_0 = \num{8640}$ blocks, i.e.\ \SI{3}{\day} at the tooling's \num{2880}
blocks per day. Reasoning: the window is the exposure period after a compromise the owner has not
noticed, so it should be short; but it is also the re-issuance cadence an autonomous agent must
survive, so an over-short window converts routine key rotation into a liveness dependency on the
owner being awake. Three days spans a weekend, which is the shortest interval that does not require
weekend attention, and it bounds undetected-compromise exposure at three days of the agent's funded
balance --- a quantity the owner already controls directly. \srcinferred.}

\begin{remark}[\FluxEdge{} scoped tokens fold into the certificate scheme --- settled]
\label{rem:fold-scoped-tokens}
\FluxEdge{}'s scoped-token model is to be \textbf{folded into this certificate scheme} rather than
kept as a separate API-layer authority \srcintent{design intake} (\planned). One
delegation object covers the platform: the same scope tuple, the same revocation semantics, and one
thing to audit rather than two systems with different trust properties to explain and keep aligned.

The cost is a real coupling and the paper states it: folding binds \FluxEdge{}'s API-layer
authorisation to the \specv{9} enforcement height (\Cref{rem:cert-location}), because the certificate
travels in a \specv{9} specification field. Until $\hstar$, \FluxEdge{} necessarily keeps its
existing tokens, so the convergence is a target state rather than a description of today. The
alternative --- decoupled schedules --- was available and was not taken, which means the two systems
must not be allowed to drift apart in the interval.
\end{remark}

\proposednote{The tooling changes, not the requirement. \Cref{rem:digest} requires $\mathcal{I}$ to
hold image digests, and the deployed tooling emits a tag; the resolution is that the tooling resolves
the tag to a digest at specification time and emits the digest. Enforcing at the node instead would
mean nodes resolving tags themselves at install time, which reintroduces exactly the mutable-reference
problem digest pinning exists to remove --- two nodes resolving the same tag at different moments can
legitimately obtain different images. \srcinferred.}

\subsection{Comparison with prior art}
\label{sec:agent-identity:compare}

\textbf{ERC-4337 session keys} \citep{erc4337} solve authority (B) directly, and the reason they can
is worth stating precisely rather than gesturing at: an EVM account is a program with persistent
state, so a validation function can read how much a session key has already spent and refuse the next
transaction on that basis. The predicate that \cref{prop:budget-balance} shows is inexpressible in
Bitcoin script is an ordinary storage read in an EVM account. \Flux's parallel-chain accounts inherit
this capability and do not currently use it; the \Flux\ L1 cannot inherit it at all. The asymmetry is
not a gap in \Flux's engineering, it is a property of the two execution models, and the paper is more
useful stating it than pretending to a uniform design across both.

\textbf{Macaroons} \citep{birgisson2014macaroons} are the closest analogue for authority (A), and the
correspondence is close enough to be worth taking seriously as a design check. A macaroon is a bearer
credential carrying attenuable caveats that a verifier checks locally, without a lookup against an
issuing service. \Cref{def:cert} is a macaroon whose caveats are specification predicates
(\cref{def:scope}), whose verification is \cref{alg:accept}, and whose revocation --- the part
macaroons conventionally handle worst, since local verification and central revocation pull against
each other --- is a UTXO spend that every verifier can already see. That is the one place where
sitting on a chain buys something a general-purpose authorisation scheme does not get for free.

\textbf{Bitcoin covenant proposals} (\texttt{OP\_CHECKTEMPLATEVERIFY}, \texttt{OP\_VAULT}) are the
primitive that would make (B) enforceable on a UTXO chain \emph{without} a new trusted party, by
letting an output constrain the transactions that spend it. \Flux\ does not have them. They are the
correct reference for what construction (c) of \cref{sec:agent-identity:spend} would cost and for what
it would buy, and they are also the reason the honest framing of row 13 in \S\ref{sec:feasibility} is ``not possible
here'' rather than ``not possible''.

\textbf{Filecoin} \citep{filecoin2017} payment channels and \textbf{Akash} \citep{akash} lease escrow
bound a client's spend by pre-funding a channel or an escrow account, which is
\cref{con:allowance} with better ergonomics and a settlement path attached. The honest framing is
that \Flux's deployable answer for spending is the same as theirs --- fund a dedicated balance, and
the balance is the bound --- and that \Flux's advantage is not there. It is in (A): collateralised
nodes, each independently validating a signed scope predicate over the specification they are about
to run, is a stronger enforcement story for deployment authority than a marketplace whose provider
accepts a lease because a broker told it to. That is the claim this section is willing to make, and
it is deliberately narrower than the one a reader might expect.

%% Drain this section's float queue before the next begins.
\FloatBarrier
